Los seis motores presentados comparten una base común de patrones de diseño y de tecnologías. El código se distribuye en varios repositorios independientes alojados en \texttt{GitHub}, cada uno con su propio entorno de ejecución, dependencias y ciclo de vida. El Anexo~\ref{cap:repositorios} recoge el mapa completo de repositorios y los enlaces de acceso a cada uno.

\subsection{Patrones de diseño}

Los motores de generación comparten varios patrones de diseño que facilitan la composición y el mantenimiento del sistema.

La separación entre generación y post-procesamiento sigue el patrón \textit{plan-then-realize}~\cite{moryossef2019stepbystep}: un componente simbólico decide la estructura de la oración y un componente neural refina el resultado. Este patrón se aplica en Hybrid1 y Hybrid2, donde la gramática o los \textit{embeddings} producen una frase estructurada y el LLM la mejora.

Hybrid3 implementa el patrón \textit{overgenerate-and-rank}~\cite{sharma2019overgenerate}: múltiples generadores producen candidatas, un filtro descarta las que no cumplen restricciones de fidelidad y una función de \textit{ranking} selecciona la mejor. Este patrón desacopla la generación de la selección, permitiendo añadir o retirar generadores sin modificar la lógica de \textit{ranking}.

Los motores que dependen de recursos costosos (modelos de \textit{embeddings} de varios GB, modelos de perplejidad) emplean inicialización perezosa: los recursos no se cargan al instanciar el motor sino en la primera llamada de generación. Esto permite instanciar múltiples \textit{pipelines} sin consumir memoria innecesariamente.

El manejo de errores sigue una estrategia de degradación progresiva: si un componente falla, el sistema recurre a un mecanismo de \textit{fallback} que garantiza siempre una salida. En las hibridaciones, si la gramática no reconoce la entrada, el LLM genera la frase desde cero; en Hybrid3, si ninguna candidata alcanza cobertura completa, todas pasan al \textit{ranking}.

\subsection{Tecnologías y dependencias}

\begin{itemize}
    \item \texttt{Python} 3.13 como lenguaje principal para los motores de generación y el \textit{pipeline} de evaluación.
    \item \texttt{Lark} para la definición y validación de la gramática formal en notación \texttt{EBNF}.
    \item \texttt{FastText}~\cite{bojanowski2017fasttext} (modelo \texttt{cc.es.300.bin}) para los \textit{embeddings} semánticos.
    \item API de OpenAI (GPT-4o, GPT-4o-mini) para el motor de LLMs y la post-edición.
    \item \texttt{spaCy} (\texttt{es\_core\_news\_sm}) para la lematización en las métricas de cobertura y fidelidad.
    \item \texttt{sacreBLEU}, \texttt{NLTK} (\textit{Natural Language Toolkit}) y \texttt{HuggingFace Transformers} para el cálculo de métricas automáticas.
    \item \texttt{React Native} y \texttt{Expo} para la aplicación móvil.
    \item \texttt{React}, \texttt{Vite} y \texttt{Supabase} para las encuestas web.
\end{itemize}

